\chapter{Multiple-Angle and Product Identities}
\label{ch:ch10-multiple-angle-product}

\section{Nothing New, Only Consequences}

The previous two chapters assembled a substantial identity toolkit: the sum
and difference formulas of Chapter~8, together with the reciprocal,
quotient, Pythagorean, cofunction, and even/odd identities of Chapter~9. At
this point trigonometry can begin to look like an ever-growing collection
of formulas with no end in sight, but that appearance is misleading. Many
of the identities still to come are consequences of a much smaller set of
principles already in hand, and the purpose of this chapter is to
demonstrate that fact explicitly. Everything derived here follows from the
sum and difference formulas of Chapter~8 together with the algebraic
identities of Chapter~9; no new geometry is introduced anywhere in this
chapter, because the geometry has already done its work.

\section{Double-Angle Formulas}

The simplest application of the sum formulas occurs when the two angles
being added are equal. Setting $A = B = \theta$ in the sine sum formula
gives
\[
\sin(2\theta) = \sin(\theta+\theta) = \sin\theta\cos\theta +
\cos\theta\sin\theta,
\]
which combines to
\[
\boxed{\sin(2\theta) = 2\sin\theta\cos\theta}.
\]
Applying the same substitution to the cosine sum formula gives
\[
\boxed{\cos(2\theta) = \cos^2\theta - \sin^2\theta}.
\]
For tangent, $\tan(2\theta) = \sin(2\theta)/\cos(2\theta)$; substituting the
two results just derived and dividing numerator and denominator by
$\cos^2\theta$ gives
\[
\boxed{\tan(2\theta) = \frac{2\tan\theta}{1-\tan^2\theta}}
\]
whenever the denominator is nonzero. None of these three formulas is a new
fact about trigonometry; each is simply the corresponding sum formula of
Chapter~8, evaluated at two equal angles.

\section{Three Forms of the Cosine Double-Angle Formula}

The cosine double-angle formula is particularly useful because the
Pythagorean identity allows it to be rewritten in two further equivalent
forms. Substituting $\sin^2\theta = 1 - \cos^2\theta$ into
$\cos(2\theta) = \cos^2\theta - \sin^2\theta$ gives
\[
\cos(2\theta) = \cos^2\theta - (1-\cos^2\theta) = 2\cos^2\theta - 1,
\]
so $\boxed{\cos(2\theta) = 2\cos^2\theta - 1}$. Substituting instead
$\cos^2\theta = 1 - \sin^2\theta$ gives
\[
\cos(2\theta) = (1-\sin^2\theta) - \sin^2\theta = 1 - 2\sin^2\theta,
\]
so $\boxed{\cos(2\theta) = 1 - 2\sin^2\theta}$. The three boxed equations
\[
\cos(2\theta) = \cos^2\theta - \sin^2\theta, \qquad
\cos(2\theta) = 2\cos^2\theta - 1, \qquad
\cos(2\theta) = 1 - 2\sin^2\theta
\]
are not three separate identities but three algebraic forms of the same
one, each convenient in different circumstances; the second and third, in
particular, will be exactly what is needed to derive the half-angle
formulas in the next section.

\begin{example}
Suppose $\sin\theta = 3/5$ and $\theta$ lies in Quadrant~I. Find
$\sin(2\theta)$. Constructing the corresponding right triangle as in
Chapter~7, the side opposite $\theta$ is $3$ and the hypotenuse is $5$, so
the remaining side is $4$ and $\cos\theta = 4/5$. Applying the double-angle
formula,
\[
\sin(2\theta) = 2\sin\theta\cos\theta = 2\left(\frac35\right)\left(\frac45\right) = \frac{24}{25}.
\]
\end{example}

\section{Half-Angle Formulas}

The half-angle formulas arise by solving the double-angle formulas in
reverse. Starting from $\cos(2\theta) = 1 - 2\sin^2\theta$ and solving for
$\sin\theta$,
\[
\sin^2\theta = \frac{1-\cos(2\theta)}{2}, \qquad
\sin\theta = \pm\sqrt{\frac{1-\cos(2\theta)}{2}}.
\]
Replacing $\theta$ by $\theta/2$ throughout gives
\[
\boxed{\sin\left(\frac{\theta}{2}\right) = \pm\sqrt{\frac{1-\cos\theta}{2}}}.
\]
The same procedure applied to $\cos(2\theta) = 2\cos^2\theta - 1$ gives
\[
\boxed{\cos\left(\frac{\theta}{2}\right) = \pm\sqrt{\frac{1+\cos\theta}{2}}},
\]
and dividing the two, $\tan(\theta/2) = \sin(\theta/2)/\cos(\theta/2)$,
gives
\[
\boxed{\tan\left(\frac{\theta}{2}\right) = \pm\sqrt{\frac{1-\cos\theta}{1+\cos\theta}}}.
\]

\section{Choosing the Correct Sign}

The half-angle formulas contain a genuine subtlety: taking a square root
introduces two possible signs, and the correct one is determined by the
quadrant of $\theta/2$, not by the quadrant of $\theta$ itself. This
distinction is essential and easy to overlook. If $\theta = 240^\circ$,
for instance, then $\theta/2 = 120^\circ$, which lies in Quadrant~II, so
$\sin(\theta/2)$ is positive while $\cos(\theta/2)$ is negative; the sign
in each half-angle formula must be chosen accordingly, based on this
quadrant of $\theta/2$ and not on the fact that $\theta$ itself lies in
Quadrant~IV. The issue plays much the same role here that the range
restriction on the inverse trigonometric functions played in Chapter~7:
a formula that looks single-valued in fact requires an extra piece of
quadrant information before it can be evaluated correctly.

\begin{algorithmproc}{How to Evaluate a Half-Angle Expression}
Identify the appropriate half-angle formula for the function requested,
and substitute the known value of $\cos\theta$ into it. Simplify the
resulting radical as far as possible, then determine the quadrant of
$\theta/2$ itself, not the quadrant of $\theta$. Choose the sign of the
radical to match that quadrant, and check the result for consistency with
the sign pattern established in Chapter~4.
\end{algorithmproc}

\begin{example}
Compute $\sin(15^\circ)$. Since $15^\circ = 30^\circ/2$, the half-angle
formula gives
\[
\sin(15^\circ) = \sqrt{\frac{1-\cos30^\circ}{2}},
\]
with the positive root chosen because $15^\circ$ lies in Quadrant~I.
Substituting $\cos30^\circ = \sqrt3/2$,
\[
\sin(15^\circ) = \sqrt{\frac{1-\frac{\sqrt3}{2}}{2}} = \sqrt{\frac{2-\sqrt3}{4}} = \frac{\sqrt{2-\sqrt3}}{2}.
\]
This agrees exactly with the value obtainable from the difference formula
of Chapter~8: the two methods are different routes to the same result, a
useful check that both chapters are internally consistent.
\end{example}

\begin{practicenow}
Compute $\cos(15^\circ)$ using the half-angle formula, and confirm that the
result matches the value obtainable from Chapter~8's difference formula.
\end{practicenow}

\section{Product-to-Sum Formulas}

The sum and difference formulas can be combined to transform a product of
trigonometric functions into a sum. Adding $\sin(A+B) = \sin A\cos B +
\cos A\sin B$ and $\sin(A-B) = \sin A\cos B - \cos A\sin B$ cancels the
mixed terms:
\[
\sin(A+B) + \sin(A-B) = 2\sin A\cos B, \qquad
\boxed{\sin A\cos B = \frac{\sin(A+B)+\sin(A-B)}{2}}.
\]
The same strategy, applied to the cosine sum and difference formulas,
gives
\[
\boxed{\cos A\cos B = \frac{\cos(A+B)+\cos(A-B)}{2}}, \qquad
\boxed{\sin A\sin B = \frac{\cos(A-B)-\cos(A+B)}{2}}.
\]
These are the product-to-sum formulas, and each is obtained purely by
adding or subtracting two instances of results already derived in
Chapter~8.

\section{Sum-to-Product Formulas}

The product-to-sum formulas can be reversed. Introducing new variables
$A = (x+y)/2$ and $B = (x-y)/2$, so that $A+B = x$ and $A-B = y$, turns the
product-to-sum formulas into statements expressing a sum as a product:
\[
\boxed{\sin x + \sin y = 2\sin\left(\frac{x+y}{2}\right)\cos\left(\frac{x-y}{2}\right)}, \qquad
\boxed{\cos x + \cos y = 2\cos\left(\frac{x+y}{2}\right)\cos\left(\frac{x-y}{2}\right)},
\]
with the analogous formulas for $\sin x - \sin y$ and $\cos x - \cos y$
following by the same substitution. The details are purely algebraic and
add nothing beyond a relabeling of the product-to-sum identities already
established.

\section{Applications}

The formulas of this chapter matter well beyond elementary trigonometry:
they simplify oscillatory expressions, support signal analysis, explain
the beat frequencies heard when two close musical pitches sound together,
appear throughout electrical engineering, and enable the calculus
techniques used to integrate products of trigonometric functions. In many
of these settings, the ability to rewrite an expression in a more
convenient form matters more than evaluating it directly, which is exactly
why a toolkit of equivalent forms, rather than a single canonical formula,
is worth having.

\section{Looking Ahead}

This chapter completes the identity toolkit begun in Chapters~8 and 9. The
book now has in hand the sum and difference formulas, the reciprocal,
quotient, Pythagorean, cofunction, and even/odd identities, and the
double-angle, half-angle, product-to-sum, and sum-to-product formulas
derived here. The next chapter finally puts this toolkit to work in a
different kind of problem. Rather than proving that two expressions are
equivalent, Chapter~11 solves equations: given a trigonometric statement
involving an unknown angle, it determines every angle for which that
statement is true, using the identities of this chapter and the last to
reduce a complicated equation to one the tools of Chapter~7 can finish.

\section{Exercises}
\begin{exercise}
Starting from the cosine sum formula, derive $\cos(2\theta) = \cos^2\theta
- \sin^2\theta$.
\end{exercise}

\begin{exercise}
Use the Pythagorean identity to derive $\cos(2\theta) = 2\cos^2\theta - 1$.
\end{exercise}

\begin{exercise}
Use the Pythagorean identity to derive $\cos(2\theta) = 1 - 2\sin^2\theta$.
\end{exercise}

\begin{exercise}
If $\sin\theta = 5/13$ and $\theta$ lies in Quadrant~I, find $\sin(2\theta)$.
\end{exercise}

\begin{exercise}
If $\cos\theta = -3/5$ and $\theta$ lies in Quadrant~II, find
$\cos(2\theta)$.
\end{exercise}

\begin{exercise}
Evaluate $\sin(22.5^\circ)$ exactly.
\end{exercise}

\begin{exercise}
Evaluate $\cos(22.5^\circ)$ exactly.
\end{exercise}

\begin{exercise}
Determine the correct sign in $\cos(\theta/2)$ when $\theta = 300^\circ$,
and explain your reasoning.
\end{exercise}

\begin{exercise}
Derive the product-to-sum formula for $\cos A\cos B$.
\end{exercise}

\begin{exercise}
Why are the double-angle and half-angle formulas best understood as
consequences of the sum formulas rather than as independent identities?
\end{exercise}
